\chapter{Phương pháp}
Trước sự ra đời của mô hình Cook-Torance vào năm 1982, các mô hình quan sát thực nghiệm như Phong \cite{phong2} và Blinn \cite{blinn_phong3} dù đạt hiệu quả về mặt tính toán, chúng thường cho ra hình ảnh trông giống nhựa (plastic-looking) và không mô phỏng chính xác tính chất vật lý của kim loại hay sự thay đổi màu sắc tại các góc nhìn nghiêng (grazing angles).

Bài báo \textit{"A Reflectance Model for Computer Graphics"} \cite{a_reflectance_model_1982} của Cook và Torrance đề xuất một phương pháp tiếp cận mới dựa trên vật lý (Physically Based Rendering - PBR). Mô hình này xem xét độc lập hai thành phần: sự phân bố năng lượng quang phổ (màu sắc) và sự phân bố định hướng của ánh sáng phản xạ. Đây là nền tảng cho các kỹ thuật kết xuất hiện đại ngày nay.

\section{Cơ sở Vật lý và Lý thuyết Vi bề mặt (Physics Basis and Microfacet Theory)}

Mô hình Cook--Torrance là một mô hình phản xạ dựa trên vật lý, được xây dựng từ quang học hình học (ray theory) và lý thuyết vi bề mặt (microfacet theory). Mục tiêu của mô hình là mô tả chính xác cách ánh sáng tương tác với các bề mặt thô nhám ở cấp độ vi mô, đồng thời vẫn đảm bảo tính khả thi trong tính toán.

\subsection{Lý thuyết Vi bề mặt (Microfacet Theory)}

Lý thuyết vi bề mặt giả định rằng một bề mặt thô nhám (rough surface) ở cấp độ vĩ mô thực chất được cấu tạo từ vô số các mặt phẳng gương siêu nhỏ gọi là các vi bề mặt (microfacets). Mỗi vi bề mặt này có pháp tuyến riêng và phản xạ ánh sáng theo định luật phản xạ gương lý tưởng.

\begin{itemize}
	\item Sự thô nhám của bề mặt được biểu diễn thông qua phân bố thống kê của các pháp tuyến vi bề mặt.
    \item Chỉ những vi bề mặt nào có vector pháp tuyến hướng đúng về phía vector nửa góc (halfway vector) giữa nguồn sáng và người quan sát mới đóng góp vào thành phần phản xạ gương nhìn thấy được.
    \begin{equation}
		\mathbf{H} = \frac{\mathbf{L} + \mathbf{V}}{||\mathbf{L} + \mathbf{V}||}
    \end{equation}
    trong đó, $\mathbf{L}$ là vector hướng nguồn sáng đơn vị và $\mathbf{V}$ là vector hướng người quan sát đơn vị.
\end{itemize}

\subsection{Các thành phần của sự phản xạ}
Dựa trên cơ chế tương tác ánh sáng với vật liệu, mô hình Cook--Torrance chia sự phản xạ thành hai thành phần chính:
\begin{enumerate}
    \item \textbf{Phản xạ gương (Specular Reflection):}\\
    Xảy ra tại bề mặt phân cách giữa không khí và vật liệu. Thành phần này mô tả ánh sáng phản xạ trực tiếp từ các vi bề mặt theo định luật phản xạ gương. Nó chịu ảnh hưởng mạnh bởi độ nhám bề mặt và chiết suất của vật liệu, đồng thời tạo ra các điểm sáng đặc trưng (specular highlights).

    \item \textbf{Phản xạ khuếch tán (Diffuse Reflection):}\\
    Xảy ra khi ánh sáng xâm nhập vào bên trong vật liệu, bị tán xạ nhiều lần và thoát ra ngoài theo nhiều hướng khác nhau. Thành phần này thường được giả định là khuếch tán đẳng hướng và ít phụ thuộc vào góc nhìn.
\end{enumerate}

\section{Mô hình Cook-Torrance (The Cook-Torrance Model)}

Mô hình Cook-Torrance định nghĩa cường độ ánh sáng phản xạ $I_r$ tới mắt người quan sát dưới tác động của nhiều nguồn sáng như sau:

\begin{equation}
I_r = I_{ia}R_a + \sum_{l} I_{il} (\mathbf{N} \cdot \mathbf{L}_l) d\omega_{il} (sR_s + dR_d)
\end{equation}

Trong đó:
\begin{itemize}
	\item $I_{ia}$ là cường độ ánh sáng môi trường (ambient light intensity).
	\item $I_{il}$ là cường độ của nguồn sáng thứ $l$.
	\item $R_a$ là hệ số phản xạ môi trường (ambient reflectance).
	\item $R_s$ là hệ số phản xạ gương hai chiều (specular bidirectional reflectance).
	\item $R_d$ là hệ số phản xạ khuếch tán hai chiều (diffuse bidirectional reflectance).
	\item $\mathbf{N}$ là vector pháp tuyến bề mặt đơn vị.
	\item $\mathbf{L}_l$ là vector hướng tới nguồn sáng đơn vị thứ $l$.
	\item $d\omega_{il}$ là góc khối của nguồn sáng, được xác định bởi:
	\begin{equation}
		d\omega_{il} = \frac{A}{r^2}
	\end{equation}
	với $A$ là diện tích bề mặt phát sáng và $r$ là khoảng cách tới điểm đang xét.
	\item $s$ và $d$ là trọng số của thành phần gương và khuếch tán, thỏa mãn $s + d = 1$.
\end{itemize}

Thành phần môi trường và khuếch tán phản xạ ánh sáng đẳng hướng, do đó $R_a$ và $R_d$ không phụ thuộc vào vị trí người quan sát. Mặt khác, thành phần phản xạ gương phản xạ ánh sáng tại một số hướng nhiều hơn phần còn lại, do đó $R_s$ phụ thuộc vào vị trí người quan sát.

Trọng tâm của mô hình Cook--Torrance nằm ở việc xây dựng hàm phản xạ gương $R_s$ dựa trên lý thuyết vi bề mặt. Thành phần này được biểu diễn dưới dạng BRDF như sau:

\begin{equation}
	R_s = \frac{F(\mathbf{V} \cdot \mathbf{H})}{\pi} \frac{D(\mathbf{H})\,G(\mathbf{N}, \mathbf{H}, \mathbf{V}, \mathbf{L})}{(\mathbf{N} \cdot \mathbf{L})(\mathbf{N} \cdot \mathbf{V})}
\end{equation}

Trong đó:

\begin{itemize}
	\item $D(\mathbf{H})$ là hàm phân bố vi bề mặt (Facet slope distribution), biểu diễn xác suất tồn tại của các vi bề mặt có pháp tuyến cùng hướng với $\mathbf{H}$. Hàm này đặc trưng cho độ nhám bề mặt.
	\item $F(\mathbf{V}, \mathbf{H})$ là hệ số Fresnel, mô tả tỉ lệ ánh sáng bị phản xạ tại bề mặt phân cách, phụ thuộc vào góc tới và chiết suất của vật liệu.
	\item $G(\mathbf{N}, \mathbf{H}, \mathbf{V}, \mathbf{L})$ là hệ số che khuất và che bóng, mô tả hiện tượng các vi bề mặt tự che khuất hoặc che bóng lẫn nhau.
	\item $(\mathbf{N} \cdot \mathbf{L})(\mathbf{N} \cdot \mathbf{V})$ đảm bảo tính bảo toàn năng lượng và chuẩn hóa BRDF.
\end{itemize}

\subsection{Hàm phân bố vi bề mặt $D$ (Normal Distribution Function)}

Hàm phân bố vi bề mặt $D$ mô tả mật độ xác suất của các vi bề mặt có pháp tuyến trùng với vector nửa góc $\mathbf{H}$, tức là mức độ “phổ biến” của các vi bề mặt đang ở trạng thái phản xạ gương lý tưởng đối với cặp hướng $(\mathbf{L}, \mathbf{V})$. Nói cách khác, $D$ quyết định cường độ của thành phần phản xạ gương thông qua việc xác định bao nhiêu vi bề mặt có thể đóng góp vào phản xạ theo hướng quan sát.

Trong mô hình Cook–Torrance, hàm phân bố Beckmann được sử dụng nhằm mô tả các bề mặt có độ nhám tuân theo phân bố Gaussian của độ dốc vi mô, dựa trên các kết quả từ lý thuyết tán xạ sóng điện từ. Hàm Beckmann được viết dưới dạng:

\begin{equation}
	D(\mathbf{H}) = \frac{1}{m^2 \cos^4 \alpha} \exp^{[-(\frac{\tan \alpha}{m})^2]}
\end{equation}

Trong đó:
\begin{itemize}
	\item $m$ là độ dốc trung bình bình phương (root mean square slope), đại diện cho mức độ nhám vi mô của bề mặt. Giá trị $m$ nhỏ tương ứng với bề mặt bóng, tập trung nhiều vi bề mặt quanh pháp tuyến trung bình. Giá trị $m$ lớn biểu thị bề mặt thô, với các vi bề mặt phân bố rộng theo nhiều hướng khác nhau.
	\item $\alpha$ là góc giữa pháp tuyến trung bình của bề mặt $\mathbf{N}$ và vector nửa góc $\mathbf{H}$, với $\cos \alpha = \mathbf{N} \cdot \mathbf{H}$.
\end{itemize}

Hàm $D$ đạt giá trị lớn nhất khi $\mathbf{H}$ gần trùng với $\mathbf{N}$, và giảm nhanh khi góc lệch tăng, phản ánh trực tiếp ảnh hưởng của độ nhám lên độ sắc nét của vùng lóe sáng gương.

\subsection{Hệ số suy hao hình học $G$ (Geometrical Attenuation Factor)}

Hệ số suy hao hình học $G$ mô hình hóa các hiện tượng che chắn hình học xảy ra trên bề mặt thô nhám, nơi các vi bề mặt có thể cản trở lẫn nhau trong quá trình truyền ánh sáng. Thành phần này đảm bảo rằng mô hình phản xạ tuân thủ các ràng buộc hình học thực tế của bề mặt vi mô.

Hai hiện tượng chính được xét đến là:
\begin{itemize}
	\item \textbf{Che bóng (shadowing):} Một vi bề mặt có hướng phản xạ phù hợp nhưng bị các vi bề mặt lân cận che khuất khỏi nguồn sáng, khiến ánh sáng tới không thể tiếp cận nó.
	\item \textbf{Che khuất (masking):} Ánh sáng đã phản xạ từ một vi bề mặt nhưng bị chặn trước khi có thể truyền đến mắt người quan sát.
\end{itemize}

Cook và Torrance đề xuất một dạng xấp xỉ cho $G$ dựa trên các ràng buộc hình học:

\begin{equation}
	G = \min \left\{
	1,\;
	\frac{2(\mathbf{N} \cdot \mathbf{H})(\mathbf{N} \cdot \mathbf{V})}{\mathbf{V} \cdot \mathbf{H}},\;
	\frac{2(\mathbf{N} \cdot \mathbf{H})(\mathbf{N} \cdot \mathbf{L})}{\mathbf{V} \cdot \mathbf{H}}
	\right\}
\end{equation}

Biểu thức này đảm bảo rằng đóng góp phản xạ gương bị suy giảm khi điều kiện hình học không cho phép cả ánh sáng tới và ánh sáng phản xạ cùng tiếp cận một vi bề mặt. Khi bề mặt trở nên bóng hơn, ảnh hưởng của $G$ giảm đi, và giá trị của nó tiến gần tới 1.

\subsection{Hệ số Fresnel $F$ (Fresnel Term)}

Hệ số Fresnel $F$ mô tả tỷ lệ ánh sáng bị phản xạ tại ranh giới giữa hai môi trường, thường là không khí và vật liệu, so với phần ánh sáng bị khúc xạ và hấp thụ vào bên trong vật liệu. Đây là thành phần phụ thuộc mạnh vào góc tới và bản chất quang học của vật liệu.

Theo các phương trình Fresnel trong quang học cổ điển, hệ số phản xạ tăng khi góc tới tăng và tiến gần đến $90^\circ$ (grazing angle), và tiến tới 1 khi ánh sáng chiếu gần song song với bề mặt. Điều này giải thích vì sao ngay cả các vật liệu kém phản xạ ở góc vuông cũng trở nên phản xạ mạnh ở các góc xiên.

Trong thực hành đồ họa máy tính, $F$ thường được xấp xỉ bằng công thức Schlick:

\begin{equation}
	F(\theta) = F_0 + (1 - F_0)(1 - \cos \theta)^5
\end{equation}

Trong đó:
\begin{itemize}
	\item $F_0$ là hệ số phản xạ tại góc tới vuông góc, đặc trưng cho tính chất vật liệu. Đối với vật liệu điện môi, $F_0$ thường nhỏ và gần như không đổi theo bước sóng. Đối với kim loại, $F_0$ lớn và phụ thuộc mạnh vào bước sóng, tạo nên màu sắc đặc trưng của phản xạ gương.
	\item $\theta$ là góc giữa vector quan sát $\mathbf{V}$ và vector nửa góc $\mathbf{H}$, với $\cos \theta = \mathbf{V} \cdot \mathbf{H}$.
\end{itemize}

Sự kết hợp của ba thành phần $D$, $G$ và $F$ tạo nên lõi vật lý của mô hình Cook–Torrance, cho phép mô phỏng phản xạ gương một cách nhất quán với các định luật quang học và là nền tảng cho các mô hình PBR hiện đại.

\section{Quy trình Kết xuất và Màu sắc (Rendering Pipeline and Color)}

Một đóng góp lớn của Cook và Torrance là phương pháp mô phỏng màu sắc dựa trên quang phổ (spectral energy distribution).

\subsection{Sự phụ thuộc vào bước sóng}
Khác với mô hình Phong gán một màu cố định cho phản xạ gương (thường là màu trắng cho nhựa hoặc màu vật liệu cho kim loại), Cook-Torrance tính toán $F$ tại từng bước sóng khả kiến.
\begin{itemize}
    \item Với \textbf{Kim loại (Metals):} Phản xạ chủ yếu diễn ra ở bề mặt. Hệ số $F$ biến đổi mạnh theo bước sóng tại góc tới thẳng đứng (tạo nên màu đặc trưng của kim loại như vàng hay đồng).
    \item Với \textbf{Phi kim/Nhựa (Plastics/Dielectrics):} $F$ gần như không đổi theo bước sóng ở góc thẳng đứng (thường là màu trắng/xám). Màu sắc của vật thể chủ yếu đến từ thành phần khuếch tán $R_d$ (do các hạt sắc tố bên dưới bề mặt).
\end{itemize}

%\subsection{Chuyển đổi sang không gian màu RGB}
%Để hiển thị trên màn hình, phân bố năng lượng quang phổ $E(\lambda)$ của ánh sáng phản xạ phải được chuyển đổi sang không gian màu XYZ (sử dụng các hàm phối màu - matching functions) và sau đó sang RGB.

\section{Ứng dụng và Kết quả (Applications and Results)}

Mô hình đã được áp dụng để mô phỏng sự khác biệt giữa kim loại và nhựa, điều mà các mô hình trước đó thường thất bại.

\begin{itemize}
    \item \textbf{Mô phỏng Nhựa (Plastics):} Sử dụng thành phần gương màu trắng (do $F$ phẳng theo bước sóng) và thành phần khuếch tán có màu. Điều này tái tạo chính xác hiệu ứng "bóng loáng" của nhựa màu.
    \item \textbf{Mô phỏng Kim loại (Metals):} Loại bỏ thành phần khuếch tán ($d \approx 0$). Thành phần gương mang màu sắc của vật liệu (ví dụ: Đồng) và thay đổi màu nhẹ khi góc nhìn chuyển sang góc nghiêng (do sự thay đổi của $F$).
\end{itemize}

%\section{Ưu điểm và Hạn chế}
%
%\subsection{Ưu điểm}
%\begin{enumerate}
%    \item \textbf{Tính chân thực vật lý:} Mô hình giải thích được hành vi ánh sáng dựa trên tính chất vật liệu thực tế (chỉ số khúc xạ $n$, hệ số tắt dần $k$).
%    \item \textbf{Thay đổi màu sắc theo góc (Color Shift):} Tái hiện chính xác hiện tượng Fresnel, làm cho kim loại trông "thực" hơn tại các cạnh viền.
%    \item \textbf{Linh hoạt:} Có thể mô phỏng đa dạng vật liệu từ nhựa bóng, cao su nhám đến kim loại rỉ sét bằng cách điều chỉnh tham số $m, s, d$.
%\end{enumerate}
%
%\subsection{Hạn chế}
%\begin{enumerate}
%    \item \textbf{Chi phí tính toán:} Tại thời điểm năm 1982, việc tính toán $F$ theo bước sóng và tích phân quang phổ rất tốn kém tài nguyên.
%    \item \textbf{Phức tạp:} Yêu cầu dữ liệu đầu vào là các chỉ số vật lý ($n, k$) hoặc quang phổ phản xạ thực nghiệm, không phải lúc nào cũng có sẵn.
%    \item \textbf{Giả định đơn giản hóa:} Mô hình vẫn giả định $R_a$ (ambient) là hằng số và bỏ qua tương tác phản xạ giữa các vật thể (inter-object reflection) ngoại trừ qua thành phần môi trường.
%\end{enumerate}

%\section{Kết luận (Conclusion)}
%
%Mô hình phản xạ Cook-Torrance đã đặt nền móng cho kỷ nguyên Kết xuất dựa trên vật lý (PBR). Bằng cách kết hợp lý thuyết vi bề mặt, phương trình Fresnel và dữ liệu quang phổ, mô hình đã chứng minh rằng việc mô phỏng chính xác các hiện tượng vật lý là chìa khóa để đạt được hình ảnh đồ họa chân thực. Các nguyên lý được trình bày trong bài báo này—đặc biệt là sự phân tách giữa thành phần gương/khuếch tán dựa trên cấu trúc vật liệu và việc sử dụng hàm phân bố $D$—vẫn là thành phần cốt lõi trong các công cụ kết xuất hiện đại (như Unreal Engine, Unity hay Blender).
